\section{Dinámica de Oscilaciones Electromecánicas}
% ===========================================================

\subsection{Dinámica de Oscilaciones Electromecánicas}
\subsubsection{Ecuaciones de Swing del Generador Síncrono}

La dinámica de cada generador $i$ en un sistema de $n_g$ máquinas ($i=1,\ldots,n_g$) se describe mediante las \textbf{ecuaciones de swing}:

\begin{block}{Ecuaciones de Swing }
\begin{align}
    \dot{\delta}_i  &= \omega_i - \omega_s \label{eq:swing1} \\[4pt]
    \dot{\omega}_i  &= \frac{\omega_s}{2H_i}
                       \Bigl(P_{m,i} - P_{e,i}(\boldsymbol{\delta}) - D_i(\omega_i - \omega_s)\Bigr)
                       \label{eq:swing2}
\end{align}
\end{block}

Para generadores síncronos modeloados de forma clásica (voltaje constante a partir de una reactancia transitoria). 

% -----------------------------------------------------------

Parámetros: 

\begin{itemize}
    \item $\delta_i$: ángulo del rotor (rad);
    \item $\omega_i$: velocidad angular del rotor (rad/s);
    \item $H_i$: constante de inercia (s); 
    \item $D_i$: constante de amortiguamiento (s);
    \item $P_{m,i}$: potencia mecánica; 
    \item $P_{e,i}(\boldsymbol{\delta})$: potencia eléctrica — función no lineal de los ángulos de rotor $\boldsymbol{\delta} = [\delta_1, \ldots, \delta_{n_g}]$;
    \item $\omega_s = 2\pi f_s$: velocidad síncrona de referencia
\end{itemize}

% -----------------------------------------------------------

% -----------------------------------------------------------

 \subsection{Dinámica de Oscilaciones Electromecánicas}
 \subsubsection{Linealización alrededor del Punto de Operación}

 Para analizar la estabilidad de pequeña señal, se linealiza el sistema alrededor del punto de equilibrio
 $(\boldsymbol{\delta}^0, \boldsymbol{\omega}^0)$, en régimen permanente se cumple ($\dot{\delta}_i = 0$, $\dot{\omega}_i = 0$): 
 \begin{equation}
    P_{m,i} = P_{e,i}(\boldsymbol{\delta}^0), \qquad \omega_i^0 = \omega_s  
 \end{equation}

Se definen las variables de estado de perturbación (desviaciones):
 \begin{equation}
    \Delta\boldsymbol{\delta} = \boldsymbol{\delta} - \boldsymbol{\delta}^0, \qquad
    \Delta\boldsymbol{\omega} = \boldsymbol{\omega} - \boldsymbol{\omega}^0
 \end{equation}

% -----------------------------------------------------------

    Linealizando la primera ecuacion de osilación:
    \begin{equation}
        \Delta\dot{\boldsymbol{\delta}} = \Delta\boldsymbol{\omega}
    \end{equation}

 Expandiendo la serie de Taylor de primero orden de $P_{e,i}(\boldsymbol{\delta})$ alrededor de $\boldsymbol{\delta}^0$ y usando la condicion de equilibrio se linealiza la segunda ecuacion de osilación:
 \begin{equation}
    \Delta\dot{\boldsymbol{\omega}}_i = \frac{\omega_s}{2H_i}\Bigl(-\sum_{j=1}^{n_g} \frac{\partial P_{e,i}}{\partial \delta_j}\big|_0 \Delta\delta_j - D_i \Delta\omega_i\Bigr)
  \end{equation}
    

% -----------------------------------------------------------

% -----------------------------------------------------------

Definimos los vectores de desviaciones: 
\begin{equation}
    \Delta\boldsymbol{\delta} = [\Delta\delta_1, \ldots, \Delta\delta_{n_g}]^\top, \qquad
    \Delta\boldsymbol{\omega} = [\Delta\omega_1, \ldots, \Delta\omega_{n_g}]^\top
\end{equation}
 las matrices diagonales:
 \begin{equation}
    \mathbf{M} = \mathrm{diag}\left(\frac{2H_i}{\omega_s}\right), \qquad
    \mathbf{K_D} = \mathrm{diag}(D_i)
 \end{equation}
 La matriz de sincronización con la red (o matriz de coeficientes de par sincronizante) representa el cambio en las potencias eléctricas para un cambio en el ángulo en el punto de equilibrio: 
\begin{equation}
    \mathbf{K}_s = \left[\frac{\partial P_{e,i}}{\partial \delta_j}\big|_{\delta = \delta^0}\right]_{i,j=1}^{n_g}
\end{equation}

% -----------------------------------------------------------

% -----------------------------------------------------------

Con estas definiciones y las ecuaciones linealizadas podemos definir la ecuacion en forma matricial: 

       \begin{equation}
    \begin{bmatrix} \Delta\dot{\boldsymbol{\delta}} \\ \Delta\dot{\boldsymbol{\omega}} \end{bmatrix}
    =
    \underbrace{
    \begin{bmatrix}
        \mathbf{0}  & \mathbf{I} \\
        -\mathbf{M}^{-1}\mathbf{K}_s & -\mathbf{M}^{-1}\mathbf{K_D}
    \end{bmatrix}
    }_{\mathbf{A} \in \mathbb{R}^{2n_g \times 2n_g}}
    \begin{bmatrix} \Delta\boldsymbol{\delta} \\ \Delta\boldsymbol{\omega} \end{bmatrix}
    \label{eq:linear_sys}
\end{equation}
Para un sistema autónomo ($u=0$) se tiene una solución de la forma: 
\begin{equation}
    \begin{bmatrix} \Delta\boldsymbol{\delta}(t) \\ \Delta\boldsymbol{\omega}(t) \end{bmatrix}
    = e^{\mathbf{A}t}\,
    \begin{bmatrix} \Delta\boldsymbol{\delta}(0) \\ \Delta\boldsymbol{\omega}(0) \end{bmatrix}
    \label{eq:solution}
\end{equation}

% -----------------------------------------------------------

Suponiendo que la matriz $\mathbf{A}$ es diagonalizable (lo que ocurre en la práctica salvo casos muy degenerados), podemos escribir
\[
\mathbf{A} = \mathbf{V} \mathbf{\lambda} \mathbf{W}^\top, \quad \mathbf{W}^\top = \mathbf{V}^{-1},
\]
donde
\begin{itemize}
    \item $\mathbf{\lambda} = \text{diag}(\lambda_1, \lambda_2, \dots, \lambda_{2n_g})$ contiene los \textbf{eigenvalores}.
    \item $\mathbf{V} = [ \mathbf{v}_1 \ \mathbf{v}_2 \ \dots \ \mathbf{v}_{2n_g} ]$ es la matriz de \textbf{eigenvectores derechos} (cada $\mathbf{v}_k$ es un vector columna de tamaño $2n_g$).
    \item $\mathbf{W} = [ \mathbf{w}_1 \ \mathbf{w}_2 \ \dots \ \mathbf{w}_{2n_g} ]$ contiene los \textbf{eigenvectores izquierdos} (vectores columna), de modo que $\mathbf{w}_k^\top$ es una fila de $\mathbf{V}^{-1}$.
\end{itemize}

La respuesta libre a partir de una condición inicial $\mathbf{x}(0)$ (la perturbación) se expresa como suma de modos:
\begin{equation}
\mathbf{x}(t) = e^{\mathbf{A}t}\mathbf{x}(0) = \sum_{k=1}^{2n_g} \mathbf{v}_k \left( \mathbf{w}_k^\top \mathbf{x}(0) \right) e^{\lambda_k t}.
\end{equation}

Esta ecuación es fundamental: muestra que el estado en cualquier instante es la superposición de los modos, cada uno con una ``intensidad'' $\alpha_k = \mathbf{w}_k^\top \mathbf{x}(0)$ (un escalar complejo si el modo es complejo; para modos reales sería real).

La ecuación de salida:
\[
y(t) = \mathbf{c}^\top \mathbf{x}(t) = \sum_{k=1}^{2n_g} \left( \mathbf{c}^\top \mathbf{v}_k \right) \left( \mathbf{w}_k^\top \mathbf{x}(0) \right) e^{\lambda_k t}.
\]

Definimos el \textbf{residuo complejo} asociado al modo $k$ como
\begin{equation}
\boxed{R_k = \left( \mathbf{c}^\top \mathbf{v}_k \right) \left( \mathbf{w}_k^\top \mathbf{x}(0) \right)}.
\end{equation}

Con esto,
\[
y(t) = \sum_{k=1}^{2n_g} R_k e^{\lambda_k t}.
\]

    Dado que los eigenvalores aparecen en pares complejos conjugados ($\lambda_{k+1} = \lambda_k^*$) y los correspondientes $R_k$ también son conjugados, podemos agruparlos para escribir la típica sinusoide amortiguada:
\[
y(t) = \sum_{\text{modos reales}} R_k e^{\sigma_k t} + \sum_{\text{pares conjugados}} 2 |R_k| e^{\sigma_k t} \cos(\omega_k t + \angle R_k).
\]

En el contexto de identificación, se suele denotar $R_{ik}$ para el canal $i$ (aquí $R_k$ para un solo canal).

Cada par conjugado $\lambda_k = \sigma_k \pm j\omega_k$ genera en el tiempo
una componente oscilatoria:
\[
y_k(t) = 2|R_k| e^{\sigma_k t} \cos(\omega_k t + \phi_k).
\]
\begin{itemize}
    \item $\omega_k$: frecuencia angular de la oscilación [rad/s].
    \item Frecuencia en \textbf{Hz}: $f_k = \dfrac{\omega_k}{2\pi}$.
    \item El decaimiento está gobernado por $\sigma_k$; para cuantificarlo
          sin unidades se usa la \textbf{razón de amortiguamiento} $\zeta_k$,
          que se define igualando el polinomio característico del modo con
          el de un oscilador de segundo orden:
          \[
          s^2 + 2\zeta_k\omega_{n,k} s + \omega_{n,k}^2 = 0,
          \]
          donde $\omega_{n,k} = \sqrt{\sigma_k^2+\omega_k^2}$.
          
\end{itemize}

% -----------------------------------------------------------

\subsection{Dinámica de Oscilaciones Electromecánicas}
\subsubsection{Eigenvalores y Modos del Sistema}

%a respuesta del sistema linealizado está gobernada por los \textbf{eigenvalores} de $\mathbf{A}$:

\begin{equation}
    \lambda_k = \sigma_k \pm j\omega_k, \qquad k = 1,\ldots,n_g
    \label{eq:eigenvalue}
\end{equation}

Los parámetros modales asociados a cada par complejo conjugado son:

\begin{block}{Frecuencia de Oscilación}
\begin{equation}
    f_k = \frac{\omega_k}{2\pi} \quad [\text{Hz}]
    \label{eq:freq}
\end{equation}
\end{block}
\begin{block}{Razón de Amortiguamiento}
\begin{equation}
    \zeta_k = \frac{-\sigma_k}{\sqrt{\sigma_k^2 + \omega_k^2}}
    \label{eq:damp}
\end{equation}
\end{block}

\begin{alertblock}{Condición de Estabilidad}
    El sistema es estable si y solo si:
    \begin{equation}
        \sigma_k < 0 \quad \forall\, k
        \label{eq:stability}
    \end{equation}
    Convencionalmente se requiere $\zeta_k \geq 5\%$.
\end{alertblock}

% -----------------------------------------------------------

\subsection{Dinámica de Oscilaciones Electromecánicas}
\subsubsection{Respuesta Temporal: Modelo de Exponenciales Amortiguadas}

La respuesta medida (velocidad angular, ángulo, potencia) de cada señal $i$ tras una perturbación puede expresarse como:

\begin{equation}
    y_i(t) = \sum_{k=1}^{n} A_{ik}\, e^{\sigma_k t}\cos(\omega_k t + \phi_{ik})
           = \mathrm{Re}\!\left\{\sum_{k=1}^{n} R_{ik}\, e^{\lambda_k t}\right\}
    \label{eq:response}
\end{equation}

donde $R_{ik} = A_{ik}e^{j\phi_{ik}}$ es el \textbf{residuo complejo} del modo $k$ en la señal $i$.

\begin{itemize}
    \item $n$: número de modos excitados por la perturbación
    \item $A_{ik}$: amplitud modal (magnitud del residuo)
    \item $\phi_{ik}$: fase modal — define la \textbf{forma modal} $\bm{\phi}_k$
\end{itemize}

    \begin{exampleblock}{Objetivo de la estimación modal}
    Recuperar $\{\lambda_k, R_{ik}\}_{k=1}^n$ directamente desde las mediciones $y_i(t)$,\textbf{ sin requerir el modelo} $\mathbf{A}$.
\end{exampleblock}

% ===========================================================
\section{Formulación Discreta y Predicción Lineal}
% ===========================================================

\subsection{Formulación Discreta y Predicción Lineal}
\subsubsection{Del Tiempo Continuo al Dominio Muestreado}

Muestreando cada señal $x_i(t)$ con periodo $\Delta t$, definimos
$x_i[n] = x_i(n\Delta t)$, para $n = 0,1,\dots,N-1$. 
Bajo el modelo de suma de exponenciales amortiguadas,
\begin{equation}
    x_i[n] = \sum_{k=1}^{M} a_{i,k}\, z_k^{\,n}, \qquad
    z_k = e^{\lambda_k \Delta t} = e^{(-\delta_k + j\omega_k)\Delta t}.
    \label{eq:discrete}
\end{equation}

Los polos en tiempo discreto $z_k$ contienen toda la información modal:

\begin{block}{Mapeo Continuo $\leftrightarrow$ Discreto}
\begin{align}
    z_k &= e^{(-\delta_k + j\omega_k)\Delta t}, \\[4pt]
    \lambda_k = -\delta_k + j\omega_k &= \frac{\ln z_k}{\Delta t}.
    \label{eq:mapping}
\end{align}
\end{block}

Esta relación permite recuperar el amortiguamiento $\delta_k$ y la frecuencia
$\omega_k$ a partir de las raíces $z_k$ del polinomio de predicción.

\subsection{Formulación Discreta y Predicción Lineal}
\subsubsection{Del Tiempo Continuo al Dominio Muestreado}

Muestreando $x_i(t)$ con periodo $\Delta t$, la muestra $x_i[n] = x_i(n\Delta t)$ satisface:

\begin{equation}
    x_i[n] = \sum_{k=1}^{M} a_{i,k}\, z_k^n, \qquad z_k = e^{\lambda_k \Delta t}
    \label{eq:discrete}
\end{equation}

La secuencia \eqref{eq:discrete} satisface una \textbf{recurrencia lineal de orden $L$}:

\begin{equation}
    x_i[n] + \sum_{\ell=1}^{L} b_\ell\, x_i[n-\ell] = 0, \qquad n \geq L
    \label{eq:ar}
\end{equation}

donde los coeficientes $\{b_\ell\}$ son las componentes del \textbf{polinomio de predicción} cuyas raíces son exactamente $\{z_k\}$.

% -----------------------------------------------------------

\subsection{Formulación Discreta y Predicción Lineal}
\subsubsection{Extracción de Parámetros Modales}

Una vez obtenidos los coeficientes $\{b_\ell\}_{\ell=1}^{L}$, los polos discretos se obtienen resolviendo el polinomio característico de orden $L$:

\begin{equation}
    P(z) = z^{L} + b_1 z^{L-1} + \cdots + b_L = 0
    \quad \Rightarrow \quad \{z_k\}_{k=1}^{L}
    \label{eq:poly}
\end{equation}

De estas $L$ raíces, solo las $M$ asociadas a la señal corresponden a modos electromecánicos.

La relación con los polos continuos $s_k = -\delta_k + j\omega_k$ está dada por $z_k = e^{s_k \Delta t}$, lo que permite extraer los parámetros:

\begin{align}
    \delta_k &= -\frac{1}{2\Delta t}\,\ln\!\bigl([\Re(z_k)]^2 + [\Im(z_k)]^2\bigr) \label{eq:delta} \\
    \omega_k &= \frac{1}{\Delta t}\,\arg(z_k) \label{eq:omega}
\end{align}

 Equivalentemente, si se introduce $\lambda_k = \frac{\ln z_k}{\Delta t} = -\delta_k + j\omega_k$, entonces
 $\delta_k = -\Re(\lambda_k)$ y $\omega_k = \Im(\lambda_k)$.

    \subsection{Formulación Discreta y Predicción Lineal}
    \subsubsection{Condición de estabilidad y parámetros finales}

    \begin{alertblock}{Condición de estabilidad en tiempo discreto}
        El sistema es estable si y solo si todos los polos $z_k$ están \textbf{dentro} del círculo unitario:
        \begin{equation}
            |z_k| < 1 \quad \forall\, k
            \label{eq:stability_z}
        \end{equation}
        lo que equivale a $\delta_k > 0$ para cada modo real.
    \end{alertblock}

    
    Una vez calculados $\delta_k$ y $\omega_k$, los parámetros modales característicos son:

    \begin{equation}
    f_k = \dfrac{\omega_k}{2\pi}\quad [Hz].
    \label{eq:freq_final}
    \end{equation}
    \begin{equation}
    \zeta_k = \dfrac{\delta_k}{\sqrt{\delta_k^2 + \omega_k^2}} \quad ( \zeta_k\% = \zeta_k \times 100).
    \label{eq:damp_final}
    \end{equation}

% ===========================================================
\section{Algoritmo Tufts-Kumaresan Multivariable (Multi-TK)}
% ===========================================================

\subsection{Algoritmo Multi-TK}
\subsubsection{Modelo de señal y predicción lineal hacia atrás}

Para $P$ señales medidas $\{x_i[n]\}_{i=1}^P$, cada una de $N$ muestras ($n=0,\dots,N-1$), el modelo de señal es
\[
x_i[n] = \sum_{k=1}^{M} a_{i,k}\, z_k^n + w_i[n], \qquad z_k = e^{(-\delta_k + j\omega_k)\Delta t}.
\]

Se elige un orden de predictor $L$ tal que $M \le L \le N-M$. Para cada canal $i$ se plantea el problema de \textbf{predicción lineal hacia atrás}:
\begin{equation}
\mathbf{A}_i\,\mathbf{b} = -\mathbf{h}_i,
\label{eq:backward_pred}
\end{equation}
donde $\mathbf{b}=[b(1),b(2),\dots,b(L)]^\top$ es el vector de coeficientes del polinomio característico.

\subsection{Algoritmo Multi-TK}

La matriz de Hankel conjugada para la $i$-ésima señal se construye como:
\begin{equation}
\mathbf{A}_i =
\begin{bmatrix}
x_i^*(1)   & x_i^*(2)   & \cdots & x_i^*(L) \\
x_i^*(2)   & x_i^*(3)   & \cdots & x_i^*(L+1) \\
\vdots     & \vdots     & \ddots & \vdots   \\
x_i^*(N-L) & x_i^*(N-L+1) & \cdots & x_i^*(N-1)
\end{bmatrix}
\in \mathbb{C}^{(N-L)\times L},
\label{eq:Ai}
\end{equation}
y el vector objetivo asociado es
\begin{equation}
\mathbf{h}_i =
\begin{bmatrix}
x_i^*(0) \\ x_i^*(1) \\ \vdots \\ x_i^*(N-L-1)
\end{bmatrix}
\in \mathbb{C}^{N-L}.
\label{eq:hi}
\end{equation}

\subsection{Algoritmo Multi-TK}
\subsubsection{Sistema global multicanal}

Se apilan verticalmente las matrices y los vectores de todos los canales:
\begin{equation}
\mathbf{A}_{\text{multi}} =
\begin{bmatrix} \mathbf{A}_1 \\ \mathbf{A}_2 \\ \vdots \\ \mathbf{A}_P \end{bmatrix}
\in \mathbb{C}^{P(N-L)\times L},
\qquad
\mathbf{h}_{\text{multi}} =
\begin{bmatrix} \mathbf{h}_1 \\ \mathbf{h}_2 \\ \vdots \\ \mathbf{h}_P \end{bmatrix}
\in \mathbb{C}^{P(N-L)}.
\label{eq:multi_system}
\end{equation}

El sistema completo de predicción lineal es:
\begin{equation}
\boxed{\mathbf{A}_{\text{multi}}\,\mathbf{b} = -\mathbf{h}_{\text{multi}}}.
\label{eq:multi_equation}
\end{equation}

% -----------------------------------------------------------

\subsection{Algoritmo Multi-TK}
\subsubsection{Descomposición SVD y selección del orden del modelo}

Se aplica la \textbf{descomposición en valores singulares} a $\mathbf{A}_{\text{multi}}$:
\begin{equation}
\mathbf{A}_{\text{multi}} = \mathbf{U}\,\boldsymbol{\Sigma}\,\mathbf{V}^H,
\label{eq:svd}
\end{equation}
con $\boldsymbol{\Sigma} = \operatorname{diag}(\sigma_1,\sigma_2,\dots,\sigma_r)$, $\sigma_1 \ge \sigma_2 \ge \dots \ge \sigma_r > 0$.

El número de modos $M$ se selecciona mediante el criterio de \textbf{energía acumulada}:
\begin{equation}
E_{s,M} = \frac{\displaystyle\sum_{k=1}^{M} \sigma_k^2}{\displaystyle\sum_{k=1}^{r} \sigma_k^2} \ge E_0,
\label{eq:energy}
\end{equation}
donde típicamente $E_0 = 97\%$ (o $99.9\%$ según el nivel de ruido).

% -----------------------------------------------------------

\subsection{Algoritmo Multi-TK}
\subsubsection{Truncamiento de la SVD y solución del sistema}

Se retiene el subespacio de señal de rango $M$:
\begin{equation}
\mathbf{A}_{\text{multi}} \approx \mathbf{U}_s \boldsymbol{\Sigma}_s \mathbf{V}_s^H,
\quad
\mathbf{U}_s \in \mathbb{C}^{P(N-L)\times M},\;
\boldsymbol{\Sigma}_s \in \mathbb{R}^{M\times M},\;
\mathbf{V}_s \in \mathbb{C}^{L\times M}.
\label{eq:truncsvd}
\end{equation}

La solución regularizada del sistema $\mathbf{A}_{\text{multi}}\,\mathbf{b} = -\mathbf{h}_{\text{multi}}$ es:
\begin{equation}
\boxed{\mathbf{b} = -\mathbf{V}_s\,\boldsymbol{\Sigma}_s^{-1}\,\mathbf{U}_s^H\,\mathbf{h}_{\text{multi}}}.
\label{eq:solution}
\end{equation}

\begin{itemize}
\item El truncamiento elimina componentes de ruido y reduce el sesgo en la estimación de $\mathbf{b}$.
\item Orden del predictor: se recomienda $L \approx \tfrac{3N}{4}$.
\end{itemize}

% ----------------------------------------------------------

\subsection{Algoritmo Multi‑TK}
\subsubsection{Solución del Sistema de Predicción Lineal}

Dado el sistema global $\mathbf{A}_{\text{multi}}\mathbf{b} = -\mathbf{h}_{\text{multi}}$,
el vector de coeficientes de predicción $\mathbf{b}\in\mathbb{C}^{L}$ se obtiene minimizando:
\begin{equation}
    \min_{\mathbf{b}}\; \|\mathbf{A}_{\text{multi}}\mathbf{b} + \mathbf{h}_{\text{multi}}\|_2^2
    \label{eq:ls}
\end{equation}

Tras la SVD truncada a $M$ componentes ($\mathbf{A}_{\text{multi}} \approx \mathbf{U}_s \boldsymbol{\Sigma}_s \mathbf{V}_s^H$):
\begin{align}
    \mathbf{w} &= \boldsymbol{\Sigma}_s^{-1}\,\mathbf{U}_s^H\,\mathbf{h}_{\text{multi}}
                 \label{eq:weights}\\[4pt]
    \mathbf{b} &= -\mathbf{V}_s\,\mathbf{w}
                 \label{eq:bcoeffs}
\end{align}

La solución retiene solo los $M$ valores singulares dominantes, lo que reduce la influencia del ruido.

\subsection{Algoritmo Multi‑TK}

Los coeficientes $\mathbf{b}=[b_1,\dots,b_L]^\top$ definen el polinomio de  predicción:
\begin{equation}
    P(z) = z^L + b_1 z^{L-1} + b_2 z^{L-2} + \cdots + b_L = 0
    \label{eq:pred_poly}
\end{equation}

\begin{exampleblock}{Raíces y modos}
    Las $L$ raíces $\{z_k\}$ de $P(z)$ son los \textbf{polos discretos} del sistema.
    De ellas, solo $M$ corresponden a la dinámica real; el resto son polos
    espurios que se filtran posteriormente mediante el criterio de energía.
\end{exampleblock}

\subsection{Algoritmo Multi‑TK}
\subsubsection{Órdenes del Predictor y del Modelo}

En el método Tufts‑Kumaresan se distingue entre:
\begin{itemize}
    \item \textbf{Orden del predictor} $L$: número de columnas de $\mathbf{A}_i$ y
          grado del polinomio $P(z)$. Se elige \emph{a priori} tal que
          $M \le L \le N-M$.
    \item \textbf{Número de modos} $M$: cantidad de componentes significativas
          (valores singulares conservados). Se determina \emph{automáticamente}
          mediante el criterio de energía acumulada (Ec.~\ref{eq:energy}).
\end{itemize}

Elegir $L > M$ introduce grados de libertad extra que permiten modelar el
ruido sin afectar las estimaciones de los verdaderos modos; las raíces
espurias se descartan en la etapa de filtrado.

\subsection{Algoritmo Multi‑TK}
\subsubsection{Resumen del Flujo Computacional}

\begin{enumerate}
    \item \textbf{Pre‑procesamiento:} Filtrado pasa‑banda $[f_{\min}, f_{\max}]$ y diezmado.
    \item \textbf{Apilamiento:} Construcción de $\mathbf{A}_{\text{multi}}$ y $\mathbf{h}_{\text{multi}}$
          desde las $P$ señales (Ec.~\ref{eq:multi_system}).
    \item \textbf{SVD:} $\mathbf{A}_{\text{multi}} = \mathbf{U}\boldsymbol{\Sigma}\mathbf{V}^H$.
    \item \textbf{Selección de orden:} $M$ mediante energía acumulada $\ge E_0$ (Ec.~\ref{eq:energy}).
    \item \textbf{Coeficientes:} Solución LS regularizada $\mathbf{b}$ (Ec.~\ref{eq:weights}--\ref{eq:bcoeffs}).
    \item \textbf{Polos discretos:} Raíces de $P(z) = z^L + b_1 z^{L-1} + \cdots + b_L$ (Ec.~\ref{eq:pred_poly}).
    \item \textbf{Parámetros modales:} Mapeo $z_k \to \lambda_k$ y cálculo de $\delta_k, \omega_k, f_k, \zeta_k$
          (Ec.~\ref{eq:delta}--\ref{eq:damp_final}.)
\end{enumerate}

%
%\begin{alertblock}{Filtrado de raíces}
 %   Solo se conservan raíces con $|z_k^{-1}| < 1$ (polos dentro del círculo unitario),
  %  lo que equivale a $\delta_k > 0$ (estabilidad en tiempo continuo).
%\end{alertblock}

% -----------------------------------------------------------

% ===========================================================
\section{Patrones Dinámicos: Formas Modales y Factores de Participación}
% ===========================================================

\subsection{Patrones Dinámicos}
\subsubsection{Estimación de amplitudes complejas (sistema de Vandermonde)}

Conocidos los polos  $\{z_k\}_{k=1}^{M}$, para cada señal $i$ ($i=1,\dots,P$)
se resuelve el sistema lineal que relaciona las muestras con las amplitudes complejas
$a_{i,k}$:

\begin{equation}
    \underbrace{
    \begin{bmatrix}
        z_1^0 & z_2^0 & \cdots & z_M^0 \\
        z_1^1 & z_2^1 & \cdots & z_M^1 \\
        \vdots &      & \ddots & \vdots \\
        z_1^{N-1} & z_2^{N-1} & \cdots & z_M^{N-1}
    \end{bmatrix}
    }_{\boldsymbol{\Psi}\,\in\,\mathbb{C}^{N \times M}}
    \begin{bmatrix} a_{i,1} \\ a_{i,2} \\ \vdots \\ a_{i,M} \end{bmatrix}
    = \mathbf{x}_i
    \label{eq:vander_art}
\end{equation}

    
La solución de mínimos cuadrados (pseudoinversa de $\boldsymbol{\Psi}$) es:

\begin{equation}
    \mathbf{a}_i = \boldsymbol{\Psi}^+\,\mathbf{x}_i
    = \bigl(\boldsymbol{\Psi}^H\boldsymbol{\Psi}\bigr)^{-1}\boldsymbol{\Psi}^H\mathbf{x}_i
    \label{eq:residues_art}
\end{equation}

De cada elemento $a_{i,k}$ se obtienen la amplitud y la fase:
\[
|a_{i,k}|,\qquad \phi_{i,k} = \arg(a_{i,k}).
\]

\subsection{Patrones Dinámicos}
\subsubsection{Forma Modal (\textit{mode shape})}

La \textbf{forma modal} del modo $k$ describe la distribución espacial de la oscilación
a través de las $P$ señales. Se construye a partir de las amplitudes complejas:

\begin{equation}
    \mathbf{s}_k =
    \begin{bmatrix}
        a_{1,k} \\
        a_{2,k} \\
        \vdots \\
        a_{P,k}
    \end{bmatrix}
    =
    \begin{bmatrix}
        |a_{1,k}|\,e^{j\phi_{1,k}} \\
        |a_{2,k}|\,e^{j\phi_{2,k}} \\
        \vdots \\
        |a_{P,k}|\,e^{j\phi_{P,k}}
    \end{bmatrix}
    \label{eq:mode_shape_art}
\end{equation}

\begin{itemize}
    \item $|\mathbf{s}_k|$ (magnitud): indica qué tan observable es el modo $k$ en cada señal.
    \item $\angle\mathbf{s}_k$ (fase): revela las relaciones de fase entre señales;
          ángulos opuestos ($\sim 180^\circ$) indican que los generadores oscilan \textbf{en contrafase}.
\end{itemize}

% -----------------------------------------------------------

\subsection{Patrones Dinámicos}
\subsubsection{Factor de Participación – definición basada en mediciones}

En el contexto de identificación a partir de señales, el \textbf{factor de participación}
se define a partir de la magnitud de las amplitudes complejas,
ya que éstas cuantifican la observabilidad relativa del modo en cada medición.

Para el modo $k$, la participación normalizada de la señal $i$ es:

\begin{equation}
    \boxed{\hat{p}_{ik} = \frac{|a_{i,k}|}{\displaystyle\sum_{j=1}^{P} |a_{j,k}|}}
    \label{eq:pf_mediciones}
\end{equation}

\begin{itemize}
    \item $\hat{p}_{ik} \in [0,1]$; un valor cercano a $1$ indica que el modo $k$
          es más observable en la señal $i$.
    \item Esta métrica permite localizar qué generadores o áreas dominan la oscilación.

\end{itemize}

% -----------------------------------------------------------

\subsection{Patrones Dinámicos}
\subsubsection{Factor de Participación – Aplicaciones}

\begin{block}{Aplicaciones}
\begin{itemize}
    \item Localización de la fuente de la oscilación.
    \item Selección óptima de la ubicación del PSS.
    \item Identificación de generadores críticos.
    \item Evaluación de observabilidad de modos en distintos puntos de la red.
\end{itemize}
\end{block}

\begin{exampleblock}{Criterio de dominancia}
    El generador $i^*$ domina el modo $k$ si:
    \begin{equation}
        i^* = \arg\max_i \hat{p}_{ik}
        \label{eq:dominancia_art}
    \end{equation}
    
    Un valor alto de $\hat{p}_{ik}$ indica que la señal $i$ es la más indicada
    para monitorear o amortiguar el modo $k$.
\end{exampleblock}

% -----------------------------------------------------------

\subsection{Patrones Dinámicos}
\subsubsection{Grupos Coherentes por Coseno Director}

Dos señales $i$ y $j$ presentan \textbf{coherencia} respecto al modo $k$
si sus amplitudes complejas son aproximadamente paralelas en el plano complejo.
El grado de coherencia se mide con el \textbf{coseno director}:

\begin{equation}
    r_{ij} =
    \frac{\displaystyle\sum_{k=1}^{M}
          \bigl(a_{i,k} - \bar{a}_i\bigr)\bigl(a_{j,k} - \bar{a}_j\bigr)^*}
         {\sqrt{\displaystyle\sum_{k=1}^{M}|a_{i,k}-\bar{a}_i|^2
                \displaystyle\sum_{k=1}^{M}|a_{j,k}-\bar{a}_j|^2}}
    \label{eq:direction_cosine_art}
\end{equation}

donde $\bar{a}_i = \frac{1}{M}\sum_{k=1}^{M} a_{i,k}$.

    \begin{itemize}
    \item $r_{ij} \approx +1$: señales completamente coherentes (oscilan en fase).
    \item $r_{ij} \approx -1$: señales en contrafase (grupos opuestos).
    \item $r_{ij} \approx 0$: comportamiento independiente respecto a los modos observados.
\end{itemize}

% -----------------------------------------------------------

\subsection{Patrones Dinámicos}
\subsubsection{Separación de Grupos Coherentes por Corte Espectral de Grafos}

Se construye un \textbf{grafo virtual no dirigido} $G = (V, E)$:

\begin{block}{Pesos de vértices y aristas}
\begin{align}
    w_{ii} &= 1 \label{eq:vertex_w}\\[2pt]
    w_{ij} &= \max(r_{ij},\; 0),\quad i\neq j \label{eq:edge_w}
\end{align}
\end{block}

\begin{block}{Laplaciano normalizado}
\begin{equation}
    \mathbf{L}_N = \mathbf{D}^{1/2}\,\mathbf{L}\,\mathbf{D}^{-1/2}
    \label{eq:laplacian}
\end{equation}
donde $L_{ij} = d_i\delta_{ij} - w_{ij}$
\end{block}

\begin{exampleblock}{Clustering espectral}
    El número de grupos $K$ se determina por el \textbf{mayor eigengap} de $\mathbf{L}_N$:
    \begin{equation}
        K = \arg\max_k\;
        (\lambda_{k+1} - \lambda_k)
        \label{eq:eigengap}
    \end{equation}
    Los grupos se obtienen aplicando $k$-means sobre los $K$ eigenvectores de $\mathbf{L}_N$.
\end{exampleblock}

% ===========================================================
\section{Refinamiento por Proyección Variable (OptVarPro)}
% ===========================================================

c
% -----------------------------------------------------------

\subsection{Refinamiento por Proyección Variable}
\subsubsection{Función de Costo y Proyección}

Sustituyendo $\mathbf{F}^*$ en el residuo, el problema se reduce a optimizar \textbf{únicamente} sobre $\bm{\lambda} \in \mathbb{C}^n$:

\begin{equation}
    \min_{\bm{\lambda}}\; J(\bm{\lambda}) =
    \frac{\bigl\|\mathbf{X}^T - \mathbf{f}(\bm{\lambda})\,\mathbf{f}^+(\bm{\lambda})\,\mathbf{X}^T\bigr\|_F^2}
         {\|\mathbf{X}\|_F^2}
    \label{eq:varpro_cost}
\end{equation}

Equivalentemente, en términos del \textbf{operador de proyección ortogonal}:

\begin{equation}
    J(\bm{\lambda}) = \frac{\bigl\|\mathbf{P}_{\mathbf{f}}^\perp\,\mathbf{X}\bigr\|_F^2}{\|\mathbf{X}\|_F^2},
    \qquad
    \mathbf{P}_{\mathbf{f}}^\perp = \mathbf{I} - \mathbf{f}\,\mathbf{f}^+
    \label{eq:projection}
\end{equation}

\begin{itemize}
    \item $J = 0$: reconstrucción perfecta
    \item El espacio de búsqueda se reduce de $2n + 2nm$ a \textbf{solo $2n$ variables}
    \item $\mathbf{F}^*$ se actualiza analíticamente en cada evaluación — sin variables de optimización adicionales
\end{itemize}

% -----------------------------------------------------------

\subsection{Refinamiento por Proyección Variable}
\subsubsection{Inicialización y Separación de Modos Dominantes}

La inicialización $\bm{\lambda}^{(0)}$ se obtiene del Multi-TK. Para separar los \textbf{modos dominantes} de los triviales (ruido) se define el \textbf{peso energético} del modo $k$:

\begin{equation}
    E_k = \|\mathbf{F}^*_k\|_2^2
    \label{eq:energy_weight}
\end{equation}

y el \textbf{peso de energía acumulada}:

\begin{equation}
    E_{s,K} = \frac{\displaystyle\sum_{k=1}^{K} E_k}{\displaystyle\sum_{k=1}^{n} E_k}
    \times 100\%
    \label{eq:cumulative_energy}
\end{equation}

\begin{block}{Criterio de selección (umbral $E_0 = 85\%$)}
    Los primeros $K^*$ modos ordenados por $E_k$ decreciente son \textbf{dominantes} si:
    \begin{equation}
        K^* = \min\bigl\{K : E_{s,K} \geq E_0\bigr\}
        \label{eq:dominant}
    \end{equation}
    Los modos con $k > K^*$ son triviales y se descartan.
\end{block}

% -----------------------------------------------------------

\subsection{Refinamiento por Proyección Variable}
\subsubsection{Restricciones y Algoritmo de Optimización}

El problema de optimización restringido es:

\begin{equation}
    \begin{aligned}
        \min_{\bm{\omega},\bm{\sigma}}\quad & J(\bm{\omega},\bm{\sigma}) \\
        \text{s.a.}\quad
        & (1-\delta_\omega)\,\omega_k^{(0)} \leq \omega_k \leq (1+\delta_\omega)\,\omega_k^{(0)} \\
        & (1-\delta_\sigma)\,\sigma_k^{(0)} \leq \sigma_k \leq (1+\delta_\sigma)\,\sigma_k^{(0)} \\
        & \sigma_k > 0,\quad \omega_k > 0, \qquad k = 1,\ldots,n
    \end{aligned}
    \label{eq:opt_problem}
\end{equation}

Las bandas $\delta_\omega, \delta_\sigma$ restringen la búsqueda alrededor de la solución TK.

\begin{block}{Solver: Trust-Region Reflective}
    Resuelve el residuo vectorial:
    \begin{equation}
        \mathbf{r}(\mathbf{p}) = \mathrm{vec}\!\left(\mathbf{P}_{\mathbf{f}}^\perp\,\mathbf{X}\right)
    \end{equation}
    Exploita la estructura de $\mathbf{r}$ para aproximar el Jacobiano eficientemente.
\end{block}

\begin{exampleblock}{Ventaja sobre Prony/ERA}
    Al minimizar $J$ directamente, OptVarPro es más robusto al ruido que métodos de predicción lineal, que minimizan implícitamente un criterio diferente.
\end{exampleblock}

% ===========================================================
\section{Resumen del Marco Teórico}
% ===========================================================

\subsection{Resumen del Marco Teórico}
\subsubsection{Integración de las Etapas}

\begin{enumerate}
    \item \textbf{Modelo físico:} Ecuaciones de swing (Ec.~\ref{eq:swing1}--\ref{eq:swing2})
          $\;\to\;$ respuesta como suma de exponenciales amortiguadas (Ec.~\ref{eq:response})
    
    \item \textbf{Discretización:} Mapeo $\lambda_k \to z_k = e^{\lambda_k T_s}$ (Ec.~\ref{eq:mapping})
          $\;\to\;$ recurrencia lineal con polinomio de predicción (Ec.~\ref{eq:ar})
    
    \item \textbf{Multi-TK:} Hankel multicanal $\to$ SVD $\to$ coeficientes LS (Ec.~\ref{eq:global}--\ref{eq:bcoeffs})
          $\;\to\;$ polos $\{z_k^{(0)}\}$ como inicialización
    
    \item \textbf{OptVarPro:} Proyección variable (Ec.~\ref{eq:varpro_cost}) con solver TR
          $\;\to\;$ polos refinados $\{z_k^*\}$
    
    \item \textbf{Patrones dinámicos:} Vandermonde $\to$ residuos $R_{ik}$ (Ec.~\ref{eq:vander})
          $\;\to\;$ formas modales (Ec.~\ref{eq:mode_shape}), factores de participación (Ec.~\ref{eq:participation}),
          grupos coherentes (Ec.~\ref{eq:direction_cosine}--\ref{eq:eigengap})
\end{enumerate}

 
    \subsection{Posible estructura del Capítulo 3}
\subsubsection{Título Capítulo 3: Casos de Estudio\\
Subtítulo: Validación del algoritmo Multi‑TK en señales sintéticas y sistemas de potencia\\}

Bloque 1 – Señales sintéticas\\
Título: Caso 1: Señales sintéticas multi-componente\\
Subtítulo: Validación con parámetros modales conocidos\\

Título: Generación de las señales de prueba\\
Subtítulo: Modos electromecánicos simulados y adición de ruido blanco\\

Título: Aplicación del Multi‑TK a las señales sintéticas\\
Subtítulo: Selección del orden mediante energía acumulada y SVD\\

Título: Resultados – Parámetros modales estimados\\
Subtítulo: Comparación de frecuencia, amortiguamiento y residuos\\

Bloque 2 – Kundur dos áreas\\
Título: Caso 2: Sistema de dos áreas de Kundur\\
Subtítulo: Análisis del modo inter‑área clásico\\

Título: Descripción del sistema y simulación transitoria\\
Subtítulo: Falla trifásica, liberación y medición de velocidad\\

Título: Identificación modal con Multi‑TK\\
Subtítulo: Estimación del modo inter‑área y modos locales\\

  Bloque 3 – IEEE 68 buses (16 generadores)\\
Título: Caso 3: Sistema IEEE de 68 barras – 16 generadores\\
Subtítulo: Validación en un sistema multi‑área (5 áreas)\\

Título: Preprocesamiento y señales transitorias\\
Subtítulo: Desviaciones de velocidad ante una falta en la línea 3‑4\\

Título: Modos electromecánicos identificados\\
Subtítulo: Frecuencias, amortiguamientos y clasificación inter‑área / local\\

Título: Patrones dinámicos: modos shape\\
Subtítulo: Diagramas polares y coherencia entre generadores\\

Título: Factores de participación\\
Subtítulo: Observabilidad modal y ubicación potencial de PSS\\

Bloque 4 – LA2030 (sistema latinoamericano)\\
Título: Caso 4: Modelo LA2030 del Sistema Interconectado Latinoamericano\\
Subtítulo: Evaluación de desempeño en un sistema real de gran escala\\

Título: Características del modelo LA2030\\
Subtítulo: Topología, áreas de control y puntos de medición\\

Título: Identificación modal en el sistema interconectado\\
Subtítulo: Modos inter‑área y oscilaciones de baja frecuencia\\

Título: Comparación con el análisis modal clásico\\
Subtítulo: Consistencia entre los modos estimados y los eigenvalores del modelo lineal\\

    Posible capitulo 4: \\

Título: Discusión y ventajas del Multi‑TK\\
Subtítulo: Robustez al ruido, precisión y eficiencia computacional\\

Título: Limitaciones y consideraciones prácticas\\
Subtítulo: Requisitos de datos, sensibilidad a la selección de orden y preprocesamiento\\

Título: Comparación con métodos tradicionales\\
Subtítulo: Prony, ERA, Eigensystem Realization Algorithm y otros\\

{Objetivo del descarte de modos espurios}
\begin{itemize}
  \item El método de Prony ajusta una señal mediante una suma de exponenciales complejas.
  \item Para un orden de modelo \(N\) se obtienen \(N\) modos (polos y residuos).
  \item Muchos modos pueden corresponder a ruido o artefactos numéricos.
  \item Se necesita un criterio para seleccionar los modos más significativos y descartar los espurios.
\end{itemize}

\section{Criterios de ordenamiento}

{Criterio 1: Amplitud del residuo}
\begin{itemize}
  \item Cada modo tiene un residuo complejo \(r_i\).
  \item La magnitud \(|r_i|\) indica la contribución instantánea del modo.
  \item Se ordenan los modos de mayor a menor \(|r_i|\).
  \item Se conservan los primeros \(\text{SUB\_N}\) modos (especificado por el usuario).
\end{itemize}
\[
\mathcal{I}_{\text{orden}} = \text{orden descendente de } |r_i|
\]

{Criterio 2: Energía del modo}
\begin{itemize}
  \item La energía de un modo se calcula como:
  \[
  E_i = |r_i|^2 \sum_{n=0}^{L-1} |p_i|^{2n}
  \]
  donde \(L\) es el número de muestras y \(p_i\) es el polo en el plano \(z\).
  \item Modos con \(|p_i|\) pequeño decaen rápido y contribuyen poca energía total.
  \item Modos persistentes (casi marginalmente estables) tienen mayor energía.
  \item Se ordenan por \(E_i\) decreciente y se eligen los primeros \(\text{SUB\_N}\).
\end{itemize}

{Comparación de criterios}
\begin{center}
\begin{tabular}{|l|p{5cm}|}
\hline
\textbf{Criterio} & \textbf{Ventaja} \\
\hline
Amplitud del residuo & Sencillo, identifica modos de alta amplitud instantánea. \\
\hline
Energía & Considera la evolución temporal; rechaza modos muy amortiguados. \\
\hline
\end{tabular}
\end{center}

En la práctica, el criterio de energía suele ser más robusto frente a ruido.

\section{Selección y reconstrucción}

{Selección de modos significativos}
\begin{enumerate}
  \item Se calculan los \(N\) modos mediante el algoritmo de Prony.
  \item Se ordenan según el criterio elegido (residuo o energía).
  \item El usuario define el número de modos a retener, \(\text{SUB\_N}\).
  \item Se toman los índices de los primeros \(\text{SUB\_N}\) modos ordenados:
  \[
  \mathcal{I}_{\text{sel}} = \mathcal{I}_{\text{orden}}(1:\text{SUB\_N})
  \]
  \item Si \(\text{SUB\_N}=0\) se conservan todos los modos.
\end{enumerate}

{Reconstrucción de la señal aproximada}
La señal reconstruida usando solo los modos seleccionados es:
\[
\hat{y}(t) = \sum_{i \in \mathcal{I}_{\text{sel}}} r_i \, e^{s_i t}
\]
donde \(s_i = \ln(p_i)/\Delta t\) es el polo en el plano continuo.

Se compara con la señal original para evaluar la calidad del ajuste (error cuadrático medio, etc.).

\section{Protección de pares conjugados}

{Integridad de pares complejos conjugados}
\begin{itemize}
  \item Para señales reales, los polos y residuos aparecen en pares conjugados.
  \item Si se selecciona solo uno de ellos, la reconstrucción sería compleja (no real).
  \item La función \texttt{check\_SUB\_N} (en \texttt{performprony.m}) evita romper pares.
\end{itemize}

{Condición de protección}
Dados dos modos consecutivos en el ordenamiento, se verifica:
\[
\begin{aligned}
&\bigl||r_k|-|r_{k+1}|\bigr| < \varepsilon_1,\\
&\left|\frac{1}{\tau_k}-\frac{1}{\tau_{k+1}}\right| < \varepsilon_2,\\
&|\omega_k-\omega_{k+1}| < \varepsilon_3,
\end{aligned}
\]
con \(\varepsilon_1=0.01\), \(\varepsilon_2=0.001\), \(\varepsilon_3=0.0001\).

Si se cumplen, se incrementa \(\text{SUB\_N}\) para incluir ambos modos (el par conjugado).

\section{Resumen del algoritmo}

{Pasos del algoritmo de descarte}
\begin{enumerate}
  \item Obtener residuos \(r_i\) y polos \(p_i\) con el método de Prony (orden \(N\)).
  \item Calcular energía \(E_i\) para cada modo (si se usa ese criterio).
  \item Ordenar modos según el criterio elegido (descendente).
  \item Seleccionar los primeros \(\text{SUB\_N}\) modos.
  \item Verificar y corregir posible ruptura de pares conjugados.
  \item Reconstruir la señal aproximada solo con los modos seleccionados.
\end{enumerate}
Este procedimiento descarta modos de baja amplitud o baja energía, considerados espurios.

{Observaciones finales}
\begin{itemize}
  \item El método es simple pero efectivo en la práctica.
  \item El criterio de energía suele proporcionar mejor separación entre modos físicos y espurios.
  \item La protección de pares conjugados evita resultados no realistas.
  \item El usuario puede ajustar \(\text{SUB\_N}\) para controlar la complejidad del modelo.
\end{itemize}

%--- CIERRE ---
\ThankYouFrame

\end{document}